% 由 final_document/recreate.md 同步；編譯：cd final_document && xelatex recreate.tex
\documentclass[11pt]{article}
\usepackage[a4paper,margin=2cm]{geometry}
\usepackage{fontspec}
\usepackage{xeCJK}
\usepackage{booktabs}
\usepackage{enumitem}

\setmainfont{Times New Roman}
\IfFontExistsTF{Songti TC}{
  \setCJKmainfont{Songti TC}
  \setCJKsansfont{Heiti TC}
}{
  \setCJKmainfont{FandolSong}
  \setCJKsansfont{FandolHei}
}

\setlist{nosep,leftmargin=1.8em}
\pagestyle{empty}

\newcommand{\term}[1]{\textbf{#1}}

\begin{document}

{\LARGE\bfseries 戰術復刻評分說明\par}
\vspace{0.8em}

\section*{1.\ 總分結構}

本關共 \term{6 題戰術}。每題分為兩項子分：

\vspace{0.4em}
\begin{tabular}{@{}lll@{}}
\toprule
\textbf{項目} & \textbf{每題滿分} & \textbf{評分內容} \\
\midrule
畫跑位 & 15 & 於戰術板標示各次接球之角色、位置及傳球方式 \\
跑戰術 & 30 & 於場地實際執行戰術 \\
\bottomrule
\end{tabular}

\vspace{0.4em}
\noindent
\textbf{合計：270 分}（45 分／題 $\times$ 6 題）。

\section*{2.\ 評分要素：「點」與「邊」}

戰術復刻僅記錄各次接球，不須標示全員跑位。

\vspace{0.4em}
\begin{tabular}{@{}p{2.2cm}p{4.2cm}p{6.8cm}@{}}
\toprule
\textbf{要素} & \textbf{定義} & \textbf{範例} \\
\midrule
點 & 接球者及其接球位置 & 前腰（AM）於 \term{右肋} 接球 \\
邊 & 該拍之前之傳球方式 & \term{傳球}、\term{直塞}、\term{傳中}、\term{倒三角傳球} 等 \\
結果（若有） & 該拍終點狀態 & \term{射正}、\term{得分} \\
\bottomrule
\end{tabular}

\vspace{0.6em}
\noindent\textbf{戰術描述格式：}

\begin{quote}
邊後衛將球 \term{傳球} 給前腰，前腰於 \term{右肋} 接球。前腰將球 \term{傳球} 給邊後衛，邊後衛於 \term{右路底線} 接球。邊後衛 \term{傳中} 給前鋒，前鋒於 \term{禁區中央} \term{得分}。
\end{quote}

\section*{3.\ 位置判定}

畫跑位與跑戰術均依 \term{20 個戰術區（zone 1--20）} 對應場地格子判定（見戰術板）。

\vspace{0.4em}
\begin{tabular}{@{}ll@{}}
\toprule
\textbf{項目} & \textbf{判定標準} \\
\midrule
畫跑位 & 需將人標於接球 zone 內；須標示接球角色及第幾拍時接球 \\
跑戰術 & 觸球瞬間 \term{球心} 須位於該 zone 對應格子內 \\
\bottomrule
\end{tabular}

\section*{4.\ 場上防守球員}

\begin{itemize}
  \item 場上配置約 \term{9 名} 靜止防守球員。
  \item 防守球員 \term{全程不移動}；其連線構成防線。
\end{itemize}

\section*{5.\ 畫跑位（15 分／題）}

\subsection*{作答要求}

\begin{enumerate}
  \item 依戰術區所載 \term{戰術文字描述} 作答。
  \item 於戰術板依序標示 \term{每一次接球}：
  \begin{itemize}
    \item \textbf{接球者}（須標示角色，如 FB、AM、ST）
    \item \textbf{接球位置}（球落於該 zone 格子內）
    \item \textbf{該拍前之傳球方式}（邊；第一拍可僅標示接球者與位置）
    \item 第幾拍時接球
  \end{itemize}
\end{enumerate}

\noindent
一題畫跑位全對得 15 分。

\section*{6.\ 跑戰術（30 分／題）}

跑戰術之戰術內容與畫跑位相同，改於場地執行，包含以下規定：

\begin{enumerate}
  \item 接球者需在正確拍數與 zone 觸球。
  \item 戰術語言的所有接球及射門必須於 10 秒內完成。
  \item 傳球者不得離開自己的 zone 傳給下一位。
  \item 撞倒防守員
  \item 球離開足球場區
\end{enumerate}

\noindent
若違反上述規定則此次嘗試失敗，共 3 次機會可以嘗試。

\end{document}
